\subsection{Apreciación Global}

El resto de esta Especificación de Requerimientos de Software se organiza en cinco capítulos principales:

\begin{itemize}
    \item \textbf{Capítulo 2 -- Descripción Global:} presenta el contexto general del producto, sus interfaces externas (sistema, usuario, software y comunicaciones), restricciones de memoria y hardware, funciones principales, perfil del usuario, modelo del dominio, suposiciones, dependencias y la matriz de distribución de requerimientos.
    \item \textbf{Capítulo 3 -- Requerimientos Específicos:} detalla de forma formal y verificable los requerimientos de interfaces externas, características funcionales del software, requerimientos de desempeño, restricciones de diseño, atributos de calidad (no funcionales) y requerimientos de la base de datos.
    \item \textbf{Capítulo 4 -- Proceso de Ingeniería de Requerimientos:} describe el enfoque metodológico basado en ISO/IEC/IEEE 29148:2018 y Disciplined Agile Delivery (DAD) para el análisis de misión, elicitación, especificación, verificación, validación y gestión de cambios.
    \item \textbf{Capítulo 5 -- Proceso de Verificación:} define la estrategia de comprobación de cumplimiento mediante revisiones e inspecciones y pruebas basadas en los estándares ISO/IEC/IEEE 12207 y 29119.
    \item \textbf{Capítulo 6 -- Anexos:} reúne material complementario y documentación de soporte asociada a la especificación.
\end{itemize}